\section{Splošno}
\s{Kitajski izrek ob ostankih}\\
Naj bo \(0 \leq x < N = n_1n_2 \cdots n_k\) in 
\begin{align*}
    x &= a_1 \mod n_1 \\
    x &= a_2 \mod n_2 \\
    &\vdots \\
    x &= a_k \mod n_k.
\end{align*}
Obstaja natanko en \(x \in \set{0, 1, ..., N - 1}\), ki hkrati zadošča vsem kongruencam. Izračunamo ga lahko po formuli
\[
    x = \sum_{i=1}^{k} a_i M_i N_i \mod N,
\]
kjer je \(N_i = N / n_i\) in \(M_i N_i = 1 \mod n_i\).\\
\textbf{Trd.} Če \(n = pq,\ p,q \in \mathbb{P}\), potem \((\Z_n, +, \cdot) \approx (\Z_p, +, \cdot) \times (\Z_q, +, \cdot)\).
\begin{itemize}
    \item Formula za izračun pogojne verjetnosti:
    \[
    P(A \mid B) = \frac{P(B \mid A) \cdot P(A)}{P(B)}, \quad P(B) > 0
    \]
    \item Formula za inverz matrike (2x2):
    \[
    A =
    \begin{pmatrix}
    a & b \\
    c & d
    \end{pmatrix}, \quad
    A^{-1} = \frac{1}{ad - bc}
    \begin{pmatrix}
    d & -b \\
    -c & a
    \end{pmatrix}
    \]

    \item Formula popolne verjetnosti:
    \[
    P(A) = \sum_{i=1}^{n} P(A \mid B_i)\, P(B_i)
    \]

    \item Formula za razcep polinoma:
    \[
    1 - x^l = (1 - x)(1 + x + x^2 + \cdots + x^{l-1})
    \]

    \item Mali Fermatov izrek:
    \[
    a^{p-1} \equiv 1 \pmod{p}, \quad \text{če } p \text{ praštevilo in } p \nmid a
    \]
\end{itemize}